%=====================================================================
\chapter{Statistical Inference Done Properly}\label{ch:inference}
%=====================================================================
\locator{4}{Part~4 --- Measurement, Data and Analysis}

\begin{outcomes}
By the end of this chapter you will be able to:
\begin{tight}
  \item \textbf{State} what a p-value is, and identify the common
    misinterpretations of it.
  \item \textbf{Report} a result as an effect size with a confidence interval
    rather than a significance verdict.
  \item \textbf{Check} the assumptions a test actually makes, and choose a
    robust alternative when they fail.
  \item \textbf{Correct} for multiple comparisons, choosing between family-wise
    and false-discovery control.
  \item \textbf{Test} for equivalence when your claim is that there is no
    meaningful difference.
\end{tight}
\end{outcomes}

\begin{prereq}
\chref{sampling} for power and effect sizes --- inference without power is
uninterpretable. \chref{measurement} for scale types, which constrain the
tests available.
\end{prereq}

\begin{keyterms}
null hypothesis significance testing $\bullet$ p-value $\bullet$ confidence
interval $\bullet$ effect size $\bullet$ assumption checking $\bullet$
normality of residuals $\bullet$ non-parametric test $\bullet$ bootstrap
$\bullet$ family-wise error rate $\bullet$ false discovery rate $\bullet$
Bonferroni $\bullet$ Holm $\bullet$ Benjamini--Hochberg $\bullet$ equivalence
testing $\bullet$ TOST $\bullet$ garden of forking paths
\end{keyterms}

\section{What a p-Value Is}

\begin{definitionbox}[The p-value]
The probability of obtaining data \emph{at least as extreme} as what you
observed, \textbf{if the null hypothesis were true} and all other assumptions of
the model held.

\medskip
That final clause is routinely dropped and is doing real work. A small p-value
indicates that the data are unusual under the \emph{entire} model --- null
hypothesis plus independence, plus the distributional assumptions, plus the
assumption that this test was the one you planned to run. Any of those can be
what is wrong.
\end{definitionbox}

\begin{paperbox}[The two documents to read on this]
Ronald L. Wasserstein and Nicole A. Lazar, \emph{The ASA Statement on p-Values:
Context, Process, and Purpose}, The American Statistician 70(2),
2016~\cite{wasserstein2016asa}.\oa{}

\medskip
Sander Greenland and colleagues, \emph{Statistical tests, P values, confidence
intervals, and power: a guide to misinterpretations}, European Journal of
Epidemiology 31, 2016~\cite{greenland2016statistical}.\oa{}

\medskip
\textbf{Why they matter.} The first is the American Statistical Association's
official position --- the professional body of statisticians stating that the
practice its own members taught for decades is widely misused. The second
enumerates twenty-five specific misinterpretations, with corrections.

\medskip
\textbf{What to extract.} From the ASA, the six principles. From Greenland,
find the misinterpretations you personally hold; most researchers find two or
three, and finding them is the point.
\end{paperbox}

\begin{center}
\small
\begin{longtable}{@{}L{6.4cm}L{7.9cm}@{}}
\toprule
\textbf{The misinterpretation} & \textbf{The correction} \\
\midrule
\endfirsthead
\toprule
\textbf{The misinterpretation} & \textbf{The correction} \\
\midrule
\endhead
\bottomrule
\endfoot
``$p$ is the probability the null hypothesis is true''
  & It is a probability computed \emph{assuming} the null is true. It says
    nothing directly about the probability of a hypothesis \\
``$p > 0.05$ means there is no effect''
  & It means this study did not detect one. With low power, that is expected
    even for real effects \\
``$p < 0.05$ means the effect is important''
  & Significance is not magnitude. With $N = 10{,}000$, a trivial effect is
    significant \\
``$p = 0.04$ and $p = 0.06$ are meaningfully different''
  & They are almost identical evidence. The threshold is a convention, not a
    boundary in nature \\
``A significant result will replicate''
  & The probability of replication depends on the true effect and the
    replication's power, not on $p$ \\
``$1 - p$ is the probability the alternative is true''
  & No. This is the same error as the first, restated \\
\end{longtable}
\end{center}

\begin{alertbox}[The ASA's six principles, compressed]
\begin{tightnum}
  \item P-values can indicate how incompatible the data are with a specified
    model.
  \item P-values do not measure the probability that the hypothesis is true, or
    that the data were produced by chance alone.
  \item Scientific conclusions should not be based only on whether a p-value
    passes a threshold.
  \item Proper inference requires full reporting and transparency --- every
    analysis run, not the one that worked.
  \item A p-value does not measure the size of an effect or the importance of a
    result.
  \item By itself, a p-value does not provide a good measure of evidence.
\end{tightnum}
\end{alertbox}

\section{Report the Estimate, Not the Verdict}

\begin{center}
\small
\begin{tabular}{@{}L{6.6cm}L{7.7cm}@{}}
\toprule
\textbf{Instead of} & \textbf{Write} \\
\midrule
``The difference was significant ($p < 0.05$).''
  & ``Checklist review found 1.3 more defects on average (95\% CI $[0.4,
    2.2]$), $d = 0.52$, $p = 0.004$.'' \\
``No significant difference was found.''
  & ``The estimated difference was 0.2 defects (95\% CI $[-1.1, 1.5]$). The
    interval includes effects large enough to matter, so this study does not
    settle the question.'' \\
``Group A outperformed group B.''
  & ``Group A found more defects, by a median difference of 2 (Cliff's $\delta
    = 0.41$, 95\% CI $[0.12, 0.63]$).'' \\
\bottomrule
\end{tabular}
\end{center}

\begin{conceptbox}[Why the interval carries more information]
A confidence interval tells you the effect size, the precision, and the
significance verdict all at once --- if the interval excludes zero, $p < 0.05$.
Nothing is lost by reporting it, and two things are gained.

\medskip
First, a null result becomes interpretable. An interval of $[-0.1, 0.1]$ says
the effect is small; an interval of $[-3.0, 3.2]$ says you learned almost
nothing. Both give $p > 0.05$, and they are entirely different findings.

\medskip
Second, your study becomes poolable. A meta-analysis (\chref{slr}) needs
estimates and their uncertainty, not verdicts. Reporting an interval is how a
small study contributes permanently to the evidence base.
\end{conceptbox}

\section{Assumptions, and What to Do When They Fail}

\begin{center}
\small
\begin{longtable}{@{}L{2.8cm}L{5.0cm}L{6.4cm}@{}}
\toprule
\textbf{Assumption} & \textbf{What it really says} & \textbf{If it fails} \\
\midrule
\endfirsthead
\toprule
\textbf{Assumption} & \textbf{What it really says} & \textbf{If it fails} \\
\midrule
\endhead
\bottomrule
\endfoot
Independence
  & Each observation contributes fresh information
  & The most serious failure and the least often checked. Model the structure:
    mixed-effects (\chref{regression}), or analyse at the unit of randomisation \\
Normality
  & The \emph{residuals} are approximately normal --- not the raw data
  & With $N$ moderate, tests are robust. Otherwise use a rank-based test,
    bootstrap the interval, or transform \\
Equal variances
  & Groups have comparable spread
  & Use Welch's t-test, which does not assume it. Welch should arguably be the
    default rather than the fallback \\
Linearity
  & The relationship has the form the model assumes
  & Plot residuals against fitted values. Add terms, or use a model that does
    not assume it \\
No influential outliers
  & No single point drives the conclusion
  & Report the analysis with and without. If they disagree, that is the finding \\
\end{longtable}
\end{center}

\begin{pitfall}[Testing for normality is usually the wrong move]
The reflex is to run Shapiro--Wilk and choose a test based on the outcome. This
is doubly unhelpful.

\medskip
With a small sample the normality test has no power, so it passes almost
anything --- exactly when non-normality would matter most. With a large sample
it detects trivial deviations and rejects, exactly when the central limit
theorem has made the parametric test robust anyway. The test's sensitivity runs
opposite to your need.

\medskip
Look at a Q--Q plot instead, and decide on the size of the departure. Better
still, use methods that do not require the assumption: rank-based tests, or
bootstrap confidence intervals, which make no distributional assumption at all
and are two lines of R.
\end{pitfall}

\begin{center}
\small
\begin{tabular}{@{}L{4.2cm}L{4.6cm}L{5.1cm}@{}}
\toprule
\textbf{Question} & \textbf{Parametric} & \textbf{Robust alternative} \\
\midrule
Two independent groups   & Welch's t-test & Mann--Whitney U; bootstrap \\
Two paired measurements  & Paired t-test  & Wilcoxon signed-rank \\
Three or more groups     & ANOVA          & Kruskal--Wallis \\
Association, two continuous & Pearson $r$ & Spearman $\rho$; Kendall $\tau$ \\
Two categorical          & Chi-square     & Fisher's exact for small counts \\
\bottomrule
\end{tabular}
\end{center}

\section{Multiple Comparisons}

Run 20 independent tests at $\alpha = 0.05$ on data with no real effects, and
you expect one significant result. This is not a subtle problem: it is
arithmetic, and it is why a paper reporting one significant finding out of many
tested comparisons has usually found nothing.

\begin{center}
\small
\begin{tabular}{@{}L{2.8cm}L{4.4cm}L{6.9cm}@{}}
\toprule
\textbf{Method} & \textbf{Controls} & \textbf{Use when} \\
\midrule
Bonferroni
  & Family-wise error rate: $\alpha/m$
  & Few comparisons, and a false positive is costly. Simple, and conservative
    to the point of destroying power when $m$ is large \\
Holm
  & Family-wise error rate, sequentially
  & Strictly more powerful than Bonferroni with the same guarantee. There is no
    reason to prefer Bonferroni over Holm~\cite{holm1979simple} \\
Benjamini--Hochberg
  & False discovery rate: expected proportion of false positives among
    rejections
  & Many comparisons, exploratory work. You accept some false positives in
    exchange for finding real effects~\cite{benjamini1995controlling} \\
None, declared
  & Nothing
  & One pre-registered primary outcome. Everything else is explicitly
    exploratory and labelled as such \\
\bottomrule
\end{tabular}
\end{center}

\begin{conceptbox}[Declaring a primary outcome is the cleanest solution]
The strongest option is not a correction at all. Nominate one primary outcome
and one primary comparison in advance (\chref{prereg}). Test it at 0.05. Report
everything else as exploratory, without significance claims.

\medskip
This is honest, it preserves power for the question you care about, and it
removes the awkward question of what constitutes a ``family'' of comparisons
--- a question with no principled answer, which is why the choice of correction
is so often argued about after the fact.
\end{conceptbox}

\begin{alertbox}[The garden of forking paths]
Multiplicity is not only about the tests you ran. It is also about the tests
you \emph{would have} run had the data looked different.

\medskip
If you would have analysed the subgroup had the main effect failed; if you would
have log-transformed had the residuals looked worse; if you would have excluded
that outlier had it gone the other way --- then your effective number of
comparisons is far larger than the number you performed, and no correction
applied to the tests you actually ran can fix it. This holds even if you ran
exactly one test and reported it honestly.

\medskip
The only real defence is to specify the analysis before seeing the data, which
is what \chref{prereg} is about, and what portfolio piece P5 asks of you.
\end{alertbox}

\section{Claiming There Is No Difference}

Failing to reject the null is not evidence for the null. To claim two things
are equivalent, you must test for equivalence.

\begin{conceptbox}[Two one-sided tests]
TOST specifies equivalence bounds --- the range within which any difference is
too small to matter, which is your SESOI from \chref{sampling} --- and tests
whether the effect is significantly \emph{inside} both
bounds~\cite{lakens2018equivalence}.

\medskip
The four possible outcomes are genuinely informative:

\begin{tight}
  \item Significant difference, not equivalent: a real, meaningful effect.
  \item Not significant, equivalent: real evidence of no meaningful effect ---
    a conclusion NHST alone can never deliver.
  \item Not significant, not equivalent: inconclusive. The study was
    underpowered, and this is the honest label for it.
  \item Significant, and equivalent: a real effect that is too small to matter.
    Common with large samples, and worth reporting plainly.
\end{tight}

\medskip
Software engineering has a genuine need for this. ``Our approach is as accurate
as the state of the art but ten times faster'' is an equivalence claim on
accuracy, and it is normally supported by a non-significant p-value, which does
not support it at all.
\end{conceptbox}

\begin{lstlisting}[language=Rlang, caption={A complete two-group analysis, reported
as an estimate. Note that the significance test is the least of what is
computed.}]
library(effsize); library(boot); library(TOSTER)

# 1. Look at the data before testing anything.
boxplot(defects ~ condition, data = d)
qqnorm(residuals(lm(defects ~ condition, data = d)))

# 2. Estimate with an interval that does not assume normality.
t.test(defects ~ condition, data = d)          # Welch by default in R
cliff.delta(defects ~ condition, data = d)     # robust effect size + CI

# 3. If the claim is "no meaningful difference", test that claim.
#    Bounds are the smallest effect of interest, not a convention.
TOSTtwo(m1 = 6.1, m2 = 5.9, sd1 = 2.4, sd2 = 2.6, n1 = 60, n2 = 60,
        low_eqbound_d = -0.40, high_eqbound_d = 0.40)

# 4. Correct only across the comparisons declared in the protocol.
p.adjust(c(p1, p2, p3), method = "holm")
\end{lstlisting}

\section{A Word on Bayesian Methods}

Bayesian analysis answers the question researchers actually want answered ---
how probable is this hypothesis given my data --- rather than the inverted one
NHST answers. It handles optional stopping gracefully, quantifies evidence for
the null rather than merely failing to reject it, and produces credible
intervals that mean what people wrongly think confidence intervals mean.

The honest caveats: it requires specifying a prior, which must be justified and
which reviewers in our field will question; conventions for reporting are less
settled; and your examiners may not know it. If your question genuinely calls
for evidence \emph{for} a null, a Bayes factor is the natural tool, and
equivalence testing is the frequentist alternative that will be more readily
accepted. Choose deliberately, and be able to defend the choice either way.

\begin{traditions}{statistical inference}
\tradEMP{Central. The discipline is to report estimates with uncertainty, check
assumptions, pre-specify comparisons, and treat significance as one piece of
information rather than a verdict.}
\tradDSR{Applies to evaluation results. A performance improvement demonstrated
without an interval is an anecdote (\chref{benchmarking}).}
\tradQUAL{Not applicable, and importing it is the category error of
\chref{qualitative}. Counting codes describes a sample; it estimates nothing.}
\tradFORM{Probability appears as an object of study rather than a tool of
inference. Randomised algorithms have provable bounds, not confidence
intervals.}
\end{traditions}

\begin{lab}[Analyse one dataset three ways]
Take a two-group dataset --- yours, or any open dataset.

\begin{tightnum}
  \item Report only ``significant / not significant''.
  \item Report the effect size with a bootstrap confidence interval.
  \item Run an equivalence test against a SESOI you justify in one sentence.
  \item Write the three conclusions side by side. Which would a practitioner
    act on?
  \item Now list every analysis decision you made --- exclusions,
    transformations, which test --- and ask how many you would have made
    differently had the data looked different. That count is your real
    multiplicity.
\end{tightnum}
\end{lab}

\begin{checkpoint}
\begin{tightnum}
  \item State precisely what $p = 0.03$ means, including the clause everyone
    omits.
  \item Two studies both report $p > 0.05$. One has CI $[-0.05, 0.05]$, the
    other $[-2.0, 2.1]$. What does each tell you?
  \item Why is testing for normality with Shapiro--Wilk usually unhelpful?
  \item You want to claim your method is as accurate as the state of the art.
    Why is a non-significant t-test insufficient, and what would suffice?
\end{tightnum}
\end{checkpoint}

\begin{chaptersummary}
\begin{tight}
  \item A p-value is the probability of data this extreme \emph{under the whole
    model}, not the probability that a hypothesis is true.
  \item Report effect sizes with confidence intervals. A null result is
    interpretable only through its interval, and only interval-reporting
    studies can be pooled.
  \item Check assumptions by looking at residuals, not by testing normality.
    Prefer Welch, rank-based methods and the bootstrap.
  \item Correct for multiplicity, preferring Holm to Bonferroni and
    Benjamini--Hochberg for exploratory work --- or declare a single primary
    outcome and label the rest exploratory.
  \item The garden of forking paths means your effective multiplicity includes
    analyses you did not run. Only pre-specification defends against it.
  \item To claim no difference, test for equivalence against a justified bound.
    Non-significance is not evidence of absence.
\end{tight}
\end{chaptersummary}

\begin{reviewq}
  \item The ASA states that scientific conclusions should not be based on
    whether a p-value passes a threshold, yet thresholds remain universal.
    Identify the institutional forces sustaining this and assess whether any
    proposed reform could overcome them.
  \item Equivalence testing requires specifying a bound, which is a judgement
    about practical importance. Critics say this makes it subjective. Is that a
    weakness, given that the alternative is a threshold nobody chose either?
  \item Software engineering data is frequently skewed and heteroscedastic, and
    parametric tests dominate anyway. What accounts for the gap between the
    field's data and its methods, and what would change it?
  \item Bayesian methods answer the question researchers want answered but
    require a prior. For a specific question in your area, specify a prior you
    could defend to an examiner, and say what you would do if they rejected it.
\end{reviewq}
